\documentclass{tubitak1001sr} % TÜBİTAK 1001 Sonuç Raporu
%
%DİKKAT DİKKAT !!!!!!!!!!!!!!! XeLaTeX --> PDF veya LuaLaTeX --> PDF yolunu izleyiniz
%
%------------------------------------------------------------------------------
\baslik{TÜBİTAK projeleri sonuç raporu LaTeX şablonu}
%       
\yurutucu{Prof. Dr. Yetkin Yürütücü}
%
\projeno{XXXFYYY}
%
\programkodu{1001, 1002, 3501, \ldots}%1001, 1002 ...
%
\arastirmacilar{Doç. Dr. Bilgin Araştırmacı \\ 
                Doç. Dr. Bilge Araştırmacı\\ Doç. Dr. İrfan Araştırmacı \\
                Dr. Öğr. Üyesi Özgün Araştırmacı
                } % \\ yerine virgül kullanılırsa araştırmacılar yanyana yazılır.
%
\danismanlar{Prof. Dr. Daniş Kamil Danışment}
%
\bursiyerler{Arş. Gör. Bilgütay Bursiyer\\ Arş. Gör. Büke Bursiyer}
%
\anahtarklm{Anahtar kelime, anahtar, kelime, anahtar kelime, anahtar, kelime,
anahtar kelime, anahtar, kelime,anahtar kelime, anahtar, kelime}
%
\keywords{Keyword, key, word, keyword, key, word, keyword, key, word,
keyword, key, word, keyword, key, word, keyword, key, word}
%
\sonucraporutarihi{MART 2023}%sonuç raporu tarihi
%------------------------------------------------------------------------------

\begin{document}
\maketitle
\input{ii_ozet}
\input{iii_abstract}
\pagenumbering{arabic}
\setcounter{page}{1}
\onehalfspacing
\input{1_giris}
\input{2_literaturozeti}
\input{3_gerecveyontem}
\input{4_bulgular}
\input{5_tartisma}
\input{6_sonucveoneriler}
%-----------------------------Kaynaklar------------------------------
\clearpage\singlespacing
\bibliographystyle{turkce_apalike}
\renewcommand{\refname}{KAYNAKLAR}
\setcounter{page}{1} 
\bibliography{kaynaklar}
%\addcontentsline{toc}{section}{KAYNAKLAR}
%-----------------------------Kaynaklar------------------------------
\end{document}